% C07 conceptual visual: hyperbolic models.
\begin{figure}[htbp]
\centering
\begin{tikzpicture}[>=Stealth,scale=0.95]
  \draw[->] (-3.1,0) -- (-0.2,0) node[right] {$\mathbb R$};
  \draw (-2.7,0) arc (180:0:1.0);
  \draw (-1.0,0) -- (-1.0,1.85);
  \node at (-1.65,2.25) {$\mathbb H$};

  \draw[->] (0.15,1.05) -- node[above] {Cayley transform} (1.9,1.05);

  \draw (3.25,1.0) circle (1.35);
  \draw (2.2,0.5) arc (205:335:1.35);
  \draw (3.25,-0.35) -- (3.25,2.35);
  \node at (3.25,2.75) {$\mathbb D$};
\end{tikzpicture}
\caption{The upper-half-plane and disk models encode the same hyperbolic geometry.
Geodesics are Euclidean lines or circles orthogonal to the ideal boundary, and the
Cayley transform carries one description to the other.}
\label{fig:vii31-c07-hyperbolic-models}
\end{figure}

\begin{table}[htbp]
\centering\small
\begin{tabular}{p{0.27\linewidth}p{0.61\linewidth}}
\toprule
Model & Metric / geodesic convention\\
\midrule
Upper half-plane & $ds^2=(dx^2+dy^2)/y^2$; geodesics are vertical lines and boundary-orthogonal semicircles.\\
Poincar\'e disk & $ds^2=4\,|dz|^2/(1-|z|^2)^2$; geodesics are boundary-orthogonal diameters or circular arcs.\\
Ideal boundary & Boundary points are not points of the metric space, although geodesics may have them as ideal endpoints.\\
Model change & Conformal maps preserve angles, not Euclidean lengths.\\
\bottomrule
\end{tabular}
\caption{Euclidean pictures differ between models; hyperbolic incidence and distance are the invariant content.}
\label{tab:vii31-c07-hyperbolic-conventions}
\end{table}
